\subsection{Problem Setup}
We consider method selection under a fixed budget parameter (budget $=6$) on Cohere command-r-plus-08-2024 runs over \texttt{openai/gsm8k} examples. GSM8K, a grade-school math benchmark, was introduced by \citet{cobbe2021training} and provides 1,319 test problems with numerical final answers that are amenable to exact-match evaluation. For each example, baseline methods produce candidate final answers. The objective is to maximize exact-match accuracy without using gold labels at runtime.

Evaluation is reported on:
\begin{itemize}
    \item final300 (seed 71),
    \item aggregate720 disjoint primary set (seeds 41, 61, 71),
    \item optional sensitivity set (+100; labeled sensitivity-only).
\end{itemize}

\subsection{Terminology}
\label{sec:terminology}

The following terms are used consistently throughout the paper.
Where helpful for readability, we use short paper-facing names and keep the corresponding implementation identifiers in Table~\ref{tab:name_mapping} for reproducibility.

\begin{description}[leftmargin=2.2cm, labelwidth=2cm, labelsep=0.2cm, style=nextline]
  \item[Candidate pool] The set of final answers produced by all available methods (frontier and external baselines) for a given example.
  \item[Frontier method] The primary, computationally intensive answer selected by the direct-reserve semantic frontier policy (\texttt{direct\_reserve\_semantic\_frontier\_v2}), which serves as the default output before any policy gate fires.
  \item[External baseline] Any of the three independently-sourced single-model methods used as reference signals: L1, S1, or TALE (defined below). External baselines provide gold-free consensus signals without access to gold labels.
  \item[L1] An external method that selects the answer with the highest support count across sampled model outputs.
  \item[S1] An external method that applies sequential budget-forcing to elicit a final answer within the allocated budget.
  \item[TALE] An external method that uses prompt-level budget cues to guide answer selection.
  \item[Logical API call] A single provider inference request counted against the budget. The budget parameter ($B = 6$) caps the maximum number of logical API calls per example across the full method stack.
  \item[Budget] The integer upper bound $B = 6$ on logical API calls per example, applied identically to all compared methods (matched-budget condition).
  \item[Matched budget] All compared methods operate under the same budget $B$, ensuring that accuracy differences are not confounded by unequal computational expenditure.
  \item[Postrun replay] Offline re-evaluation of a previously generated run artifact using a policy rule without issuing new provider API calls. Postrun replay is used to compute policy-specific accuracy from existing logs.
  \item[Gold-free runtime feature] Any input to a policy decision rule that is derived solely from model outputs, run metadata, or structural properties of the candidate pool---never from gold labels, exact-match signals, or correctness indicators.
  \item[Promotion-grade evidence] Evidence collected on held-out evaluation sets that were not used during rule development, meeting the disjointness, seed-independence, and aggregate scale criteria documented in Section~\ref{sec:experiments}.
  \item[Provider-specific evidence] Results obtained from a particular provider--model pairing (here, Cohere command-r-plus-08-2024). Claims are explicitly scoped to this setting and do not generalize beyond it without additional replication.
\end{description}
